\section{Auditing the Judge Against Independent Frameworks}
\label{sec:judgeaudit}

Every result up to this point rests on a single measurement device, which
is an LLM judge that maps a free-text response to a binary verdict. If
that judge is idiosyncratic, then the entire difficulty scale inherits
its idiosyncrasy. In this section I ask how far that problem actually
extends, by scoring the same responses with two evaluation frameworks
that I did not write.

\subsection{Design}
I sampled 400 responses from the Section~\ref{sec:phase3b} log, which is
100 frozen-bank items crossed with 4 zero-shot cells (two models at two
strictness levels), and re-scored each response with four continuous
metrics drawn from two external libraries.

\begin{itemize}[nosep,leftmargin=1.4em]
  \item \textbf{G-Eval strict} \citep{liu2023geval}, a chain-of-thought
        rubric that, like my judge, requires a section-level regulatory
        citation.
  \item \textbf{G-Eval lenient}, the same rubric with the citation
        requirement removed, scoring the conclusion and reasoning only.
  \item \textbf{RAGAS faithfulness} \citep{es2024ragas}, which measures
        how well the response is grounded in the supplied scenario
        context.
  \item \textbf{RAGAS answer correctness}, which measures agreement with
        the gold standard. I use the factual component only, with no
        embedding-similarity term.
\end{itemize}

No new solver calls were needed, because the audit reuses responses that
were already collected. All four metrics run on the same base model as my
judge (\texttt{llama-3.3-70b-versatile}). That choice is deliberate.
Holding the grader's raw capability fixed isolates the effect of the
\emph{rubric} from the effect of the \emph{model}. It also bounds what
the audit can show, and I return to that point in
Section~\ref{sec:judgeaudit-limits}.

I treat the judge's binary verdict as the reference and ask how well each
continuous metric separates \textsc{pass} from \textsc{fail}. I report
the area under the ROC curve (AUC) and the point-biserial correlation. Of
the $1{,}600$ scoring calls, 93 ($5.8\%$) failed on transient errors and
are excluded pairwise, so the RAGAS metrics have a smaller $n$.

\subsection{Results}

\begin{table}[h]
\centering
\small
\caption{External metrics against the judge's binary verdict on 400
         zero-shot responses. AUC is the probability that a randomly
         chosen \textsc{pass} response scores above a randomly chosen
         \textsc{fail} response, so $0.5$ is chance.}
\label{tab:judgeaudit}
\begin{tabular}{lccc}
\toprule
Metric & $n$ & AUC vs.\ judge verdict & $r_{pb}$ \\
\midrule
G-Eval strict (citation rubric)   & 400 & \textbf{0.838} & $+0.570$ \\
G-Eval lenient (conclusion only)  & 400 & 0.731 & $+0.381$ \\
RAGAS answer correctness          & 344 & 0.693 & $+0.315$ \\
RAGAS faithfulness (groundedness) & 363 & 0.530 & $+0.057$ \\
\bottomrule
\end{tabular}
\end{table}

Table~\ref{tab:judgeaudit} shows a clean ordering, and it is the ordering
the design predicts.

\paragraph{Convergent validity.}
The citation-requiring G-Eval rubric tracks my judge most closely, at
$\mathrm{AUC} = 0.838$. Two rubrics written independently, sharing only
the stipulation that a compliance answer must cite the governing
provision, agree on which responses pass roughly five times out of six.
Removing the citation clause from that same rubric drops the agreement to
$0.731$. That drop is direct evidence that what my judge enforces is the
citation requirement itself, rather than some incidental stylistic
preference.

\paragraph{Discriminant validity.}
RAGAS faithfulness sits essentially at chance
($\mathrm{AUC} = 0.530$, $r_{pb} = +0.06$). This result is the reason the
metric was included. A response can be perfectly grounded in the scenario
it was given, inventing no facts and contradicting nothing, and still
reach the wrong compliance conclusion. Groundedness and correctness are
different constructs, and a metric that measures the first carries almost
no information about the second. I read this as a caution about metric
selection in regulated domains rather than as a defect of RAGAS, which is
measuring what it claims to measure.

\paragraph{Difficulty is not confounded with quality.}
A leniency artifact would show up as external scores drifting with item
difficulty \emph{within} a verdict class. I regressed each metric's score
on $b$ separately within the \textsc{pass} and \textsc{fail} responses.
For G-Eval strict the slopes are near zero ($+0.015$, $p = 0.18$ within
\textsc{pass}; $-0.019$, $p = 0.03$ within \textsc{fail}). The remaining
metrics show small but significant negative slopes within \textsc{fail},
the steepest being RAGAS faithfulness at $-0.073$ ($p < 10^{-7}$), which
means that failing responses to harder items are scored slightly lower
still. So the IRT separation of $b$ from response quality is not
perfectly reproduced by the external tools, but neither is it strongly
violated.

\paragraph{Cell-level ordering.}
Averaging each metric within the four examinee cells and ranking those
means against the Rasch $\theta$ of the same cells, both G-Eval variants
reproduce the $\theta$ ordering exactly (Spearman $\rho = 1.00$) and
RAGAS faithfulness inverts it ($\rho = -0.80$). With only four cells this
is descriptive rather than inferential, since any monotone metric has a
one-in-twelve chance of perfect agreement, and I report it as a
consistency check rather than as evidence.

\subsection{What the audit does and does not establish}
\label{sec:judgeaudit-limits}
The audit rules out one specific failure mode, which is that my judge's
verdicts might be arbitrary, or that they might encode a preference
unrelated to the stated rubric. An independently written citation-aware
rubric recovers those verdicts at $\mathrm{AUC} = 0.838$, and a
groundedness metric correctly does not.

The audit does not establish that the verdicts are \emph{right}. All four
metrics run on the same base model, and all four consume the same
LLM-authored gold standards. A shared blind spot in the underlying model
family, or an error in a gold standard, is invisible to every metric in
Table~\ref{tab:judgeaudit}. Convergence among rubrics that share a
substrate is weaker evidence than convergence across substrates, and it
is weaker still than agreement with a human expert. Closing that gap is
the purpose of Section~\ref{sec:expertreview}.

So that the residual disagreement can be inspected rather than merely
acknowledged, I release the 40 strongest conflicts, which are the
responses scoring high on G-Eval strict but judged \textsc{fail} and the
converse, in \texttt{judge\_audit\_disagreements.jsonl}.
